% Vietnamese translation for OI Public Library
% Source-File: IOI中国国家候选队论文/2006/胡伟栋/胡伟栋.ppt
\documentclass[11pt]{article}
\usepackage{oipl-vn}

\title{Bản trình chiếu của Hu Weidong}
\author{Hu Weidong}
\date{2006}

\begin{document}
\maketitle

\origtitle{Hu Weidong presentation}
\sourcepath{IOI China candidate team papers / 2006 / Hu Weidong / Hu Weidong.ppt}

\section{Tình trạng nguồn}

Tệp PowerPoint cũ chỉ cung cấp một lớp văn bản hạn chế: ngày 2006.1.19, một số
trường liên hệ, các biểu thức như \(n^2\), \(1\cdot2\cdot3\), và vài chuỗi
tên trang. Phần lớn nội dung có khả năng nằm trong hình, hiệu ứng hoặc đối
tượng nhúng. Vì vậy bản dịch này chỉ ghi lại những mốc kỹ thuật có thể kiểm
chứng trực tiếp từ tệp, không tự thêm một bài toán hay thuật toán không có
trong lớp văn bản trích được.

\section{Các mốc đọc được}

Những mảnh còn đọc được gợi ý rằng bài trình bày có liên quan tới phân tích
độ lớn, tăng trưởng hàm và so sánh cách tính. Các biểu thức \(n^2\) và
\(1\cdot2\cdot3\) cho thấy slide có thể dùng ví dụ về bậc thời gian hoặc tích
lũy qua nhiều bước, nhưng văn bản trích xuất không đủ để khẳng định đề bài cụ
thể.

\section{Nội dung có thể giữ lại}

Các ký hiệu đọc được phù hợp với kiểu slide nhập môn về đánh giá chi phí:
\begin{itemize}
  \item \(n^2\) thường đại diện cho số cặp hoặc hai vòng lặp lồng nhau;
  \item \(1\cdot2\cdot3\) minh họa phép nhân liên tiếp hoặc quá trình tích
  lũy theo từng bước;
  \item các trường liên hệ và ngày tháng cho thấy đây là bộ slide trình bày,
  không phải một bài viết độc lập có phần phát biểu đầy đủ.
\end{itemize}

Do không có bản DOC tương ứng trong thư mục 2006, phần dịch không dựng thêm
chi tiết thuật toán. Nếu sau này có ảnh từng slide hoặc bản chuyển đổi đáng
tin cậy, cần thay phần này bằng bản dịch theo từng trang.

\section{Hướng đọc tiếp}

Nếu sau này chuyển được slide thành ảnh hoặc PDF, có thể đọc trực quan các
hình và khôi phục nội dung theo từng trang. Ở trạng thái hiện tại, mọi phần
trình bày sâu hơn đều cần nguồn hiển thị tốt hơn.

\end{document}
